% Part: proof-theory
% Chapter: natural-deduction
% Section: grafting

\documentclass[../../../include/open-logic-section]{subfiles}

\begin{document}

\olfileid[hi]{pt}{nat}{gra}

\olsection{\usetoken{P}{proof} की कलमबंदी}

\Log{N1c} और~\Log{N1i} में !!^{proof}, मान्यताओं~$!B$ से !!{formula}s~$!A$
के !!{proof}s हैं। मान लें स्वयं $!B$ का कोई !!a{proof} है। हम~$!B$ के
!!{proof} को~$!B$ से $!A$ के !!{proof} के साथ कैसे मिलाएँ कि~$!A$ का ऐसा
!!a{proof} मिले जो मान्यता~$!B$ का उपयोग न करे?

एक विकल्प है~\Intro{\lif} लागू करके~$!A$ के !!{proof} से मान्यता~$!B$ को
मुक्त करना। फिर~$!B \lif !A$ के फलित !!{proof} को~$!B$ के !!{proof} के साथ
मिलाकर यह मिलता है
\[
\AxiomC{$\Discharge{!B}{x}$}
\DeduceC{$!A$}
\DischargeRule{\Intro{\lif}}{x}
\UnaryInfC{$!B \lif !C$}
\AxiomC{}
\DeduceC{$!B$}
\RightLabel{\Elim{\lif}}
\BinaryInfC{$!A$}
\DisplayProof
\]
यह सीधा है, किंतु कुछ असुंदर: $\lif$ को प्रस्तुत करके तुरंत उसका विलोपन ही
क्यों करें? $!B \lif !A$ से होकर इस प्रकार के ``चक्कर'' से बचना संभव होना
चाहिए। (वास्तव में उन्हें सदा टाला जा सकता है, यही प्रसामान्यीकरण प्रमेय की
विषयवस्तु है।)

\Log{N1c} और~\Log{N1i} में इस मामले में हम इसे काफ़ी सरलता से टाल सकते हैं:
मान्यताओं~$!B^x$ को केवल~$!B$ के पूरे !!{proof} से बदल सकते हैं। अन्य
!!{proof}s की मान्यताओं पर !!{proof}s की यह ``कलमबंदी'' काफ़ी उपयोगी है,
अतः हम इसे परिभाषा के रूप में दर्ज करते हैं।

\begin{defn}
यदि $\delta_1$, खुली मान्यता~$!B^x$ से~$!A$ का कोई
\Log{N1c}-!!{proof} (या \Log{N1i}-!!{proof}) है और $\delta_2$,~$!B$ का
कोई \Log{N1c}-!!{proof} (\Log{N1i}-!!{proof}) है, तो
$\Subst{\delta_1}{\delta_2}{!B^x}$ वह परिणाम है जो~$\delta_1$ में~$!B^x$
की प्रत्येक अमुक्त उपस्थिति को~$\delta_2$ से बदलने पर मिलता है।
\end{defn}

बशर्ते $\delta_1$ और $\delta_2$ में एक ही चिह्न वाले भिन्न मान्यता
!!{formula}s न हों और~$\delta_2$ की किसी खुली मान्यता में~$\delta_1$ का कोई
आइगेनचर न हो, फलित $\Subst{\delta_1}{\delta_2}{!B^x}$ एक सही !!{proof} है
और उसकी खुली मान्यताएँ $\delta_1$ तथा~$\delta_2$ की खुली मान्यताएँ हैं।
!!{proof}s की कलमबंदी करते समय ये शर्तें सामान्यतः संतुष्ट होंगी। पहली तब
संतुष्ट होती है जब $\delta_1$ और $\delta_2$ किसी बड़े !!{proof} के उपप्रमाण
हों। यदि दूसरी संतुष्ट न हो, तो~$\delta_1$ के आइगेनचरों के नाम बदलकर उसे
संतुष्ट कर सकते हैं।

एक संगत परिभाषा \Log{N2c} और~\Log{N2i} पर लागू होती है।

\begin{defn}
यदि $\delta_1$,~$x:!B, \Gamma_1 \Sequent !A$ का कोई
\Log{N2c}-!!{proof} (या \Log{N2i}-!!{proof}) है और $\delta_2$,
~$\Gamma_2 \Sequent !B$ का कोई \Log{N1c}-!!{proof}
(\Log{N1i}-!!{proof}) है, तो $\Subst{\delta_1}{\delta_2}{x:!B}$ वह परिणाम
है जो~$\delta_1$ के प्रत्येक अभिगृहीत $x:!B \Sequent !B$ को~$\delta_2$ से
बदलकर और ऐसे अभिगृहीतों के नीचे प्रत्येक अनुक्रम के प्रसंग में $\Gamma_2$
जोड़कर मिलता है।
\end{defn}

\begin{prop}
यदि $\delta_1 \Proves x:!B, \Gamma_1 \Sequent !A$ और $\delta_2 \Proves
\Gamma_2 \Sequent !B$, तो $\Subst{\delta_1}{\delta_2}{x:!B} \Proves
\Gamma_1, \Gamma_2 \Sequent !A$, बशर्ते~$\delta_1$ का कोई आइगेनचर
~$\Gamma_2$ में न आए।
\end{prop}


\begin{proof}
  $\pheight{\delta}$ पर आगमन से।
\end{proof}

\end{document}
